%%=============================================================================
%% Literatuurevaluatie van de PoC
%%=============================================================================

\chapter{\IfLanguageName{dutch}{Literatuurevaluatie van de Proof-of-Concept}{Literature-based evaluation of the Proof-of-Concept}}
\label{ch:evaluatie}

In dit hoofdstuk wordt de Proof-of-Concept uit hoofdstuk~\ref{ch:poc} geëvalueerd in twee delen.
Het eerste deel is een \emph{vergelijkende literatuurstudie} die in deze bachelorproef is uitgevoerd: elke ontwerpkeuze in \textit{Learn COBOL} wordt afgezet tegen empirische bevindingen over gamification, programmeeronderwijs en didactiek (secties~\ref{sec:evaluatie-methode}--\ref{sec:evaluatie-didactisch}).
Het tweede deel is een \emph{uitgewerkt onderzoeksdesign} voor een pilotstudie waarin de PoC empirisch wordt vergeleken met traditionele methoden zoals een papieren cursus met dezelfde leerthema's (sectie~\ref{sec:evaluatie-pilotstudie}).
Samen leveren beide delen een onderbouwd antwoord op de centrale onderzoeksvraag (sectie~\ref{sec:onderzoeksvraag}) en de deelvragen: de literatuurevaluatie op theoretische plausibiliteit, het pilotontwerp op hoe die plausibiliteit met deelnemers gemeten kan worden.

\section{Doel en aanpak van de evaluatie}
\label{sec:evaluatie-doel}

De evaluatie heeft drie concrete doelstellingen:
\begin{itemize}
    \item \textbf{Theoretische ondersteuning:} per gamification-element en per Octalysis Core Drive nagaan of de gemaakte ontwerpkeuze door de literatuur wordt ondersteund, gerelativeerd of tegengesproken.
    \item \textbf{Antwoord op de deelvragen, met name DV5:} inschatten in hoeverre \textit{Learn COBOL} plausibel effectiever is dan traditionele methoden zoals een papieren cursus, op basis van vergelijkbaar empirisch onderzoek.
    \item \textbf{Voorbereiding van empirische validatie:} een reproduceerbaar pilotontwerp vastleggen zodat de literatuurhypotheses later met deelnemers gemeten kunnen worden.
\end{itemize}

De koppeling met de deelvragen uit sectie~\ref{sec:onderzoeksvraag} is als volgt.
DV1, DV2, DV3 en DV4 worden inhoudelijk al beantwoord in respectievelijk hoofdstuk~\ref{ch:stand-van-zaken} en hoofdstuk~\ref{ch:poc}.
DV5 wordt in dit hoofdstuk beantwoord in twee stappen: eerst via de literatuurevaluatie (meta-analyse van \textcite{Zhan2022}, verschil tussen interactieve en statische leertools uit hoofdstuk~\ref{ch:stand-van-zaken}), daarna via het pilotontwerp dat beschrijft \emph{hoe} die vergelijking empirisch zou worden uitgevoerd.

\paragraph{Evaluatiestrategie: literatuurevaluatie en pilotontwerp}
\label{subsec:evaluatie-keuze}

De oorspronkelijke onderzoeksopzet voorzag een empirische vergelijking tussen de PoC en traditionele methoden zoals een papieren cursus: twee groepen deelnemers leren elk één dag COBOL, één groep via de PoC en één groep via een papieren cursus, en hun resultaten op dezelfde kennistest worden vergeleken.
Dat volledige onderzoek kon binnen het tijdsbestek van deze bachelorproef niet worden uitgevoerd.
Twee knelpunten bepaalden die keuze:

\begin{itemize}
    \item \textbf{Rekrutering en planning.} Voldoende deelnemers vinden die aan de inclusiecriteria voldoen vraagt veel screening. De doelgroep bestaat uit junior IT'ers met Python-ervaring én zonder enige COBOL-kennis. Voor een zinvolle spreiding wordt typisch \(N \geq 12\) als minimum gehanteerd.
    \item \textbf{Materiaal en uitvoering.} Een eerlijke vergelijking vereist een vooraf gevalideerde papieren cursus, screening, gestandaardiseerde instrumenten en voor elke deelnemer twee geplande sessies plus export van applicatielogs.
\end{itemize}

In plaats van een onderbouwde pilot met slechts twee of drie deelnemers uit te voeren, is gekozen voor een tweeledige aanpak.


Literatuurevaluatie (uitgevoerd): de ontwerpkeuzes van de PoC worden theoretisch \emph{gekalibreerd} tegen bestaande empirische literatuur.
Die stap levert geen causale claim over \textit{Learn COBOL} zelf, wel een conceptueel onderbouwde uitspraak: indien de PoC ingezet wordt zoals beschreven, is het op basis van vergelijkbaar onderzoek aannemelijk dat de gamification-elementen een positief effect hebben op motivatie en leerwinst.
Pilotontwerp: sectie~\ref{sec:evaluatie-pilotstudie} beschrijft een vergelijkend experiment in uitvoerbare vorm (doelgroep, steekproef, tijdslijn, meetinstrumenten en analyse), zodat een opvolger de empirische validatie kan overnemen zonder het onderzoeksplan opnieuw te moeten ontwerpen.

\section{Methode van de literatuurevaluatie}
\label{sec:evaluatie-methode}

De evaluatie volgt een eenvoudig drie-staps-protocol per ontwerpkeuze:

\begin{enumerate}
    \item Beschrijf de ontwerpkeuze in de PoC zoals geïmplementeerd (zie hoofdstuk~\ref{ch:poc}).
    \item Identificeer een empirische claim uit de literatuur die op die keuze van toepassing is.
    \item Beoordeel de uitlijning tussen de ontwerpkeuze en de claim: bevestigt de literatuur de keuze, zet ze die in perspectief, of waarschuwt ze ertegen?
\end{enumerate}

De voornaamste bronnen voor stap~2 zijn de meta-analyse van \textcite{Zhan2022} (over gamification in programmeeronderwijs), het werk van \textcite{KevinWerbach2020} (over de PBL-triade en haar valkuilen), het Octalysis-framework van \textcite{Chou2015}, het MDA-framework van \textcite{Hunicke2004}, het overzicht van \textcite{Cepeda2006} over het spacing effect en de marktanalyse van COBOL-leermiddelen uit hoofdstuk~\ref{ch:stand-van-zaken}.


\paragraph{Wat de literatuurevaluatie niet doet.}
Ze meet géén leeruitkomsten op deelnemers en levert dus géén testresultaten of andere directe metingen op gebruikers op.
De secties~\ref{sec:evaluatie-pbl}--\ref{sec:evaluatie-didactisch} geven per ontwerpkeuze een \emph{verwacht effect}, niet een gemeten effect op gebruikers van \textit{Learn COBOL}.

\section{Evaluatie van de PBL-elementen}
\label{sec:evaluatie-pbl}

Het protocol uit sectie~\ref{sec:evaluatie-methode} wordt hieronder toegepast op vier thema's: PBL, Octalysis, MDA en het didactische ontwerp.

De PoC implementeert de drie klassieke PBL-componenten: punten (met first-try bonus), badges (zeven inhoudsbadges, plus combinatiebadges rond clean run, mijlpalen en speedrun) en een leaderboard. 
\textcite{KevinWerbach2020} benadrukken dat PBL het \emph{startpunt} van gamification is, niet het eindpunt: enkel een dunne PBL-laag op een bestaand systeem leggen (de zogenaamde \textit{PBL-fallacy}) leidt zelden tot duurzame motivatie. 
De vraag is dus niet of PBL aanwezig is, maar of de drie elementen ingebed zijn in een breder ontwerp dat onderliggende motivatoren raakt.

\paragraph{Punten}
De PoC kent als beloning voor een geslaagd level 100 punten toe en geeft 25 bonuspunten wanneer de eerste indiening correct is. Deze keuze sluit aan bij wat \textcite{KevinWerbach2020} beschrijven als ``zinvolle'' puntenstructuren: punten zijn gekoppeld aan een concrete prestatie.
De first-try bonus maakt daarnaast gebruik van de Core Drive \textit{Loss \& Avoidance} \autocite{Chou2015}: door de bonus te missen verlies je een onmiddellijk zichtbaar voordeel, wat aanzet tot zorgvuldig lezen vóór indienen. 

\paragraph{Badges}
De PoC kent zeven inhoudsbadges (één per leerthema) plus combinatiebadges. 
Hiernaast is er een pagina waar gebruikers verdiende en nog te verdienen badges kunnen bekijken. Hier kunnen ze ook zien hoe ze de badge kunnen verdienen.
Dit komt overeen met de aanbevelingen van \textcite{KevinWerbach2020} dat badges twee functies moeten combineren: een erkenning- en een ontdekkingsfunctie.

\paragraph{Leaderboard}
Dit is het meest betwiste PBL-element. \textcite{KevinWerbach2020} waarschuwen expliciet dat leaderboards in zakelijke contexten kunnen demotiveren wanneer de afstand tot de top te groot is, en in extreme gevallen zelfs ongewenst gedrag kunnen uitlokken. 
De PoC mitigeert dit risico op twee manieren: het volledige leaderboard staat op een aparte route (\texttt{/leaderboard}), en het dashboard toont een ingekorte top-drie van vrienden in plaats van de globale top-N. 

\paragraph{Tussentijdse conclusie.}
De PBL-implementatie in \textit{Learn COBOL} vermijdt de PBL-fallacy door alle drie de elementen aan een bredere ontwerplaag te koppelen (zie sectie~\ref{sec:evaluatie-octalysis} en~\ref{sec:evaluatie-mda} hieronder).
De keuzes worden door de literatuur over PBL ondersteund.

\section{Evaluatie tegen het Octalysis-framework}
\label{sec:evaluatie-octalysis}

Het Octalysis-framework van \textcite{Chou2015} biedt acht \textit{Core Drives} (CD's) als checklist om te beoordelen of een gamified systeem gebalanceerd is.
Tabel~\ref{tab:eval-octalysis} koppelt elk Core Drive aan de concrete ontwerpkeuzes in de PoC.

\begin{table}[H]
\centering
\small
\begin{enumerate}[label=\textbf{CD\arabic*}, leftmargin=*, itemsep=0.55em, topsep=0.4em]
    \item \textit{Epic Meaning \& Calling}: mainframe-strips rond \texttt{WORKING-STORAGE} en \texttt{PROCEDURE DIVISION}; narratief dat COBOL kritieke infrastructuur draagt.
    \item \textit{Development \& Accomplishment}: levels, voortgangsbalk, checklist met vinkjes, punten, badges, viering bij succes, rang en tier op het profiel.
    \item \textit{Empowerment of Creativity \& Feedback}: vrije code-invoer met onmiddellijke validatie tegen leerdoelen.
    \item \textit{Ownership \& Possession}: profielpagina met avatar, cumulatieve punten en een badge-collectie gekoppeld aan het account.
    \item \textit{Social Influence \& Relatedness}: vriendenlijst, top-drie vrienden op het dashboard, volledig leaderboard op een aparte route, profielen van vrienden.
    \item \textit{Scarcity \& Impatience}: sequentiële ontgrendeling van levels, \textit{locked} badges blijven zichtbaar maar onbereikbaar.
    \item \textit{Unpredictability \& Curiosity}: optionele tips en documentatielinks per level; variatie in badge-eisen (speedrun, mijlpalen).
    \item \textit{Loss \& Avoidance}: bewust beperkt: geen ``levens'' of accountreset; enkel een first-try bonus die verloren kan gaan en lichte sociale vergelijking via het leaderboard.
\end{enumerate}
\caption{Mapping van Octalysis Core Drives op ontwerpkeuzes in \textit{Learn COBOL}}
\label{tab:eval-octalysis}
\end{table}

\paragraph{White hat / black hat balans.}
\textcite{Chou2015} adviseert een combinatie van \textit{White Hat} motivatoren (CD1, CD2, CD3: positief en versterkend) en \textit{Black Hat} motivatoren (CD6, CD7, CD8: gebaseerd op schaarste, nieuwsgierigheid en verliesangst). 
De PoC is duidelijk White-Hat-zwaar. Core Drives 1--3 (tabel~\ref{tab:eval-octalysis}) richten zich op voortgang, feedback en betekenis, terwijl CD6--CD8 bewust beperkt zijn ingezet. 

\paragraph{Verwachte effectgrootte.}
De meta-analyse van \textcite{Zhan2022} rapporteert voor gamification in programmeeronderwijs een gestandaardiseerd gemiddeld effect van ongeveer \(SMD = 0{,}77\) op studentmotivatie en \(\approx 0{,}75\) op academische prestaties, beide medium-tot-grote effecten. 
Aangezien de PoC alle door Zhan beschreven kerncomponenten (punten, badges, levels, directe feedback, sociale vergelijking) implementeert, is het aannemelijk dat een eventuele empirische test minstens binnen dezelfde grootteorde zou vallen. 
Deze voorspelling moet wel met voorbehoud gelezen worden: de meta-analyse aggregeert sterk uiteenlopende studies en contexten, en de overdraagbaarheid naar specifiek COBOL-onderwijs is niet getest.

\section{Evaluatie via het MDA-perspectief}
\label{sec:evaluatie-mda}

Het MDA-framework \autocite{Hunicke2004} dwingt om na te denken over hoe \textit{mechanics} (regels en code) zich vertalen in \textit{dynamics} (gedrag tijdens gebruik) en uiteindelijk in \textit{aesthetics} (de emotionele beleving).

\begin{itemize}
    \item \textbf{Mechanics:} validatiecode, puntberekening, sequentiële ontgrendeling.
    \item \textbf{Dynamics:} competitie om rang, het mikken op een first-try bonus, het vergelijken van eigen voortgang met die van vrienden, herhaalbaar oefenen van een level.
    \item \textbf{Aesthetics:} \textit{Challenge} (levels overwinnen) als dominante emotie, \textit{Fantasy} via het mainframe-thema en \textit{Narrative} via de strips.
\end{itemize}

De evaluatie levert hier vooral een kwalitatieve check op: zijn de mechanics ontworpen om de gewenste aesthetics te ondersteunen?
In de PoC zorgt de combinatie van checklist, voortgangsbalk en celebration-overlay ervoor dat \textit{Challenge} en de bijhorende beloning in dezelfde momenten gebeuren. 
Dit ontwerpprincipe, mechanics laten samenkomen naar één moment, is wat \textcite{Hunicke2004} ``the right kind of fun'' noemen.

\section{Evaluatie van het didactische ontwerp}
\label{sec:evaluatie-didactisch}

Naast de gamification-laag bevat de PoC een aantal expliciet didactische keuzes die los van Octalysis of MDA staan.

\paragraph{Heuristische validatie en checklistfeedback.}
De keuze om de PoC niet te koppelen aan een echte COBOL-compiler maar te valideren via regex-patronen op leerdoelen, levert twee voordelen op: feedback is direct (geen compilatietijd) en feedback is in dezelfde \emph{taal} geformuleerd als de opdracht (``alle vier divisions aanwezig'') in plaats van compilermeldingen die voor beginners cryptisch overkomen.
Onderzoek over interactieve programmeeromgevingen toont dat het type feedback minstens even belangrijk is als de aanwezigheid ervan: feedback die expliciet leerdoel-gericht is, wordt door beginners beter benut dan generieke foutmeldingen \autocite{Zhan2022}. 
De keerzijde is een verlies aan correctheidsgarantie: de PoC kan code goedkeuren die in een echte compiler zou falen.

\paragraph{Distributed practice en spacing effect.}
De levelstructuur met herhaalbare opdrachten, badge-eisen die meerdere bezoeken vereisen (bv.\ speedrun), en het feit dat een voltooid level later opnieuw kan worden geopend, ondersteunen distributed practice. 
\textcite{Cepeda2006} beschrijven in hun overzicht van het spacing effect dat leren over meerdere kortere sessies typisch betere lange-termijn retentie oplevert dan één blokje massaal oefenen. 
Een gamified leerpad met streaks, mijlpalen en herhaalbare levels is in dat opzicht een betere container voor verspreide oefening dan een statisch PDF-document, dat doorgaans één keer wordt doorlopen.

\section{Beperkingen van de literatuurevaluatie}
\label{sec:evaluatie-beperkingen}

De literatuur-gebaseerde resultaten in secties~\ref{sec:evaluatie-pbl}--\ref{sec:evaluatie-didactisch} hebben inherente beperkingen.
Die beperkingen zijn geen afwijzing van empirisch onderzoek, maar motiveren het pilotontwerp in sectie~\ref{sec:evaluatie-pilotstudie}.

\begin{itemize}
    \item \textbf{Geen empirische geldigheid voor de PoC zelf.} De resultaten ondersteunen de plausibiliteit van de ontwerpkeuzes, maar tonen niet aan dat \textit{Learn COBOL} in de praktijk een hogere leer-effectiviteit haalt dan een statisch alternatief.
    \item \textbf{Selectiebias in de literatuur.} De geraadpleegde meta-analyses focussen vooral op moderne talen (Python, Java, Scratch, \ldots). Of de gemeten effecten zich op dezelfde manier voordoen bij een verbose, decennia-oude taal als COBOL is niet onderzocht.
    \item \textbf{Confounding van interactiviteit en gamification.} De positieve effecten in \textcite{Zhan2022} ontstaan deels door interactiviteit op zich, niet enkel door gamification.
    \item \textbf{Heuristische validatie.} De PoC valideert COBOL-code via regex-patronen, niet via een echte compiler. False positives en false negatives zijn mogelijk.
\end{itemize}

\section{Ontwerp van een pilotstudie ter validatie van de PoC}%
\label{sec:evaluatie-pilotstudie}

De literatuurevaluatie levert plausibiliteit. Een empirische meting op echte deelnemers vereist het onderzoeksdesign dat hieronder wordt uitgewerkt.
Het plan sluit aan op de oorspronkelijke onderzoeksopzet (PoC versus papieren cursus, zelfde zeven leerthema's) en op de beperkingen uit sectie~\ref{sec:evaluatie-beperkingen}.

Onderstaand plan beschrijft een beknopte pilotstudie met een beperkte steekproef (\(N = 12\) tot \(20\)).
Het doel is na te gaan of de PoC tot betere leeruitkomsten leidt dan een klassieke aanpak na één leerdag.
De vergelijking gebeurt op één enkele kennistest die beide groepen maken na hun leerdag.

\subsection{Doelstellingen}

De pilot heeft één gericht doel: nagaan of de PoC tot betere leeruitkomsten leidt dan een klassieke aanpak na één leerdag.
Concreet wordt vergeleken hoe deelnemers die één dag met de PoC werken scoren op dezelfde kennistest als deelnemers die één dag met een papieren cursus werken, met dezelfde zeven leerthema's en volgorde als in hoofdstuk~\ref{ch:poc}.

\subsection{Doelgroep, inclusiecriteria en steekproef}

\textbf{Doelgroep:} junior developers of IT-studenten met een Python-basis en geen voorafgaande formele COBOL-opleiding.

\textbf{Inclusie:} zelfgerapporteerd voldoende Python om de parallelle voorbeelden te volgen, en bereidheid om één leerdag en aansluitend de kennistest af te leggen.

\textbf{Exclusie:} deelnemers aan een Mainframe-/COBOL-specialisatie, professionele COBOL-ervaring.

\textbf{Steekproefgrootte:} We streven naar \(N = 12\) tot \(20\) deelnemers in totaal, ongeveer 6 tot 10 per groep.

\subsection{Onderzoeksopzet}

Het ontwerp is tussen-proefpersoon. Dat wil zeggen dat elke deelnemer maar één van de twee leermethoden volgt, niet beide. Zo wordt vermeden dat de eerste methode de tweede meting beïnvloedt.

De deelnemers worden willekeurig verdeeld over twee groepen:
\begin{itemize}
    \item \textbf{Groep A (klassiek):} leert COBOL gedurende één dag met traditionele methoden zoals een papieren cursus die dezelfde zeven leerthema's behandelt als in hoofdstuk~\ref{ch:poc}.
    \item \textbf{Groep B (PoC):} leert COBOL gedurende één dag met \textit{Learn COBOL}, langs dezelfde zeven leerthema's in dezelfde volgorde.
\end{itemize}

Aan het eind van de leerdag maken beide groepen dezelfde kennistest over die zeven leerthema's. De primaire vergelijking is het gemiddelde testresultaat van groep A tegenover dat van groep B.

\subsection{Tijdslijn}

Figuur~\ref{fig:pilot-tijdlijn} geeft een relatieve tijdslijn van vijf opeenvolgende weken vanaf de start van de pilotvoorbereiding.

\begin{figure}[H]
\centering
\resizebox{\linewidth}{!}{%
\begin{tikzpicture}[
    font=\small,
    tick/.style={thick},
    lbl/.style={align=left, text width=3.15cm, inner sep=2pt}
]
    \coordinate (A) at (-0.15,0);
    \coordinate (END) at (14.25,0);
    \draw[tick,-{Latex[length=2.8mm]}] (A) -- (END);
    \foreach \x in {0, 4.05, 8.1, 11.6}
        \draw (\x,-0.17) -- (\x,0.17);
    \node[lbl, below=14pt] at (0,-0.17) {\textbf{Week~1}\\[3pt]
        Papieren cursus en kennistest afronden, informed consent opstellen, instrumenten testen buiten de steekproef.};
    \node[lbl, above=14pt] at (4.05,0.17) {\textbf{Week~2}\\[3pt]
        Screening van deelnemers en willekeurige toewijzing aan groep A of groep B.};
    \node[lbl, below=14pt] at (8.1,-0.17) {\textbf{Week 3--4}\\[3pt]
        Datacollectie: één leerdag per deelnemer in de toegewezen groep, gevolgd door de kennistest.};
    \node[lbl, above=14pt] at (11.6,0.17) {\textbf{Week~5}\\[3pt]
        Analyse, vergelijking van de testresultaten tussen beide groepen en pilotrapport.};
\end{tikzpicture}%
}%
\caption{Relatieve tijdslijn van de pilotfase}%
\label{fig:pilot-tijdlijn}
\end{figure}

\subsection{Meetinstrumenten}

\paragraph{Kennistest na de leerdag.}
Aan het eind van de leerdag legt elke deelnemer dezelfde kennistest af. Dat is een korte test met vragen over de zeven leerthema's uit tabel~\ref{tab:poc-levels}. Beide groepen krijgen exact dezelfde vragen, zodat de scores rechtstreeks vergelijkbaar zijn.

De primaire uitkomstmaat is het gemiddelde aantal correct beantwoorde vragen per groep. Het verschil tussen beide groepsgemiddelden geeft een indicatie van de leer-effectiviteit van de PoC ten opzichte van de klassieke papieren cursus.